El hallazgo principal no es que una representación domine universalmente a las demás, sino que diferentes representaciones parecen capturar distintos aspectos de la tarea. TF--IDF más regresión logística ofrece el mejor ordenamiento, lo que sugiere que la señal léxica es muy informativa en este dominio. Esto es plausible porque los documentos de licitación suelen contener fórmulas administrativas repetidas, requisitos estandarizados, terminología específica y expresiones procedimentales que pueden correlacionarse fuertemente con los resultados posteriores. En este contexto, características léxicas dispersas pueden recuperar gran parte de la señal predictiva mediante un modelo lineal relativamente simple.

Esto no debe interpretarse como evidencia contraria a las representaciones basadas en LLM. Todos los modelos no triviales de esta familia permanecen claramente por encima de las líneas base y los embeddings promediados se mantienen próximos al transformer. Ello indica que los embeddings congelados de los fragmentos ya codifican información densa útil. En otras palabras, buena parte del valor del LLM puede residir en la propia representación, más que en la complejidad del mecanismo posterior de agregación.

La comparación entre el promedio y el transformer resulta especialmente informativa. Si las interacciones detalladas entre fragmentos fueran la principal fuente de capacidad predictiva, el transformer debería dominar claramente las líneas base agregadas. Eso no ocurre en los experimentos actuales. El MLP con promedio se aproxima al transformer en ordenamiento y la regresión logística con promedio incluso lo supera en exactitud balanceada. Esto sugiere que el modelado explícito de interacciones aún no proporciona una ventaja clara de ordenamiento bajo el régimen de datos y la señal de supervisión actuales.

No obstante, el transformer conserva una ventaja importante: la calibración y la estabilidad del umbral. Obtiene el mejor puntaje de Brier y muestra muy poca sensibilidad al ajuste del umbral, mientras que TF--IDF más regresión logística se beneficia considerablemente de dicho ajuste en F1. Desde una perspectiva aplicada, esto importa porque el modelo preferido para ordenar no necesariamente es el mismo que se prefiere para priorizar mediante probabilidades o para lograr un comportamiento estable en producción. La naturaleza relativamente estandarizada de los documentos de construcción vial también puede contribuir al buen desempeño de las líneas base léxicas.

Parte de la señal predictiva puede reflejar la complejidad documental, las prácticas institucionales de redacción, la estandarización de procedimientos u otras regularidades textuales latentes. Estas señales son legítimas para la predicción porque están disponibles ex ante, cuando se emite el documento. Sin embargo, no deben interpretarse como causas de los resultados de una contratación.

En conjunto, la comparación resulta informativa incluso sin un único ganador dominante. Para aplicaciones de gobierno y democracia electrónicos, este es precisamente el punto: los sistemas de monitoreo de contrataciones requieren más que un puntaje alto de ordenamiento. También necesitan probabilidades calibradas, líneas base transparentes y umbrales cuyo comportamiento pueda explicarse antes de utilizarlos para priorizar la atención humana. Los resultados muestran que los modelos léxicos son muy competitivos, que los embeddings congelados de LLM ya capturan gran parte de la señal densa útil y que la principal ventaja comparativa del transformer reside en su calibración y estabilidad ante el umbral, no en su desempeño bruto de ordenamiento.
